%% kapa: exam=1 round=36 year=2025 period=1 form=A subject=civil_law subject_ko=민법 seq=22 number=22
%% source: Q-Net 감정평가사 기출문제 (cst00309) artlSeq=5231525 file=raw/1008/5231525/2025년 제36회 감정평가사 1차시험 1교시 문제지 원본.pdf
%% conversion: kordoc 4.24.51 -> tools/gen1.py (auto); review=proofread-sonnet; flags=none
\begin{question}{22}
부동산 물권변동을 위하여 등기가 필요하지 않은 경우를 모두 고른 것은? (다툼이 있으면 판례에 따름)
\begin{framedtext}
ㄱ. 부동산 매매계약이 해제되어 소유권이 매도인에게 복귀하는 경우 \par
ㄴ. 화해조서에 의하여 부동산 소유권을 취득하는 경우 \par
ㄷ. 공유자 사이에 현물분할에 관한 조정이 성립하여 이에 따라 각 부동산의 단독소유권을 취득하는 경우
\end{framedtext}
\begin{choices}
  \item ㄱ
  \item ㄷ
  \item ㄱ, ㄴ
  \item ㄴ, ㄷ
  \item ㄱ, ㄴ, ㄷ
\end{choices}
\end{question}
\answer{1}
